%%%%%%%%%%  scan idx 449  |  folio 436  %%%%%%%%%%
Comme les deux membres de la relation à démontrer sont additifs, on peut supposer que \(x\) est la classe d'un \(\cO_Y\)-Module localement libre de type fini \(M\). L'élément \(i^*(i_*(x))\) est alors la classe du complexe \(\mathbb L i^*(\mathbb Ri_*(M))\), obtenu en prenant une résolution projective du complexe \(\mathbb Ri_*(M)=i_*(M)\), et en la tensorisant par \(\cO_Y\). Les faisceaux d'homologie de ce complexe sont donc les \(\underline{\Tor}_i^{\cO_X}(M,\cO_Y)\). D'après 2.4, on a l'isomorphisme:
\[
\underline{\Tor}_i^{\cO_X}(M,\cO_Y)
\xrightarrow{\sim}
M\otimes_{\cO_Y}\underline{\Tor}_i^{\cO_X}(\cO_Y,\cO_Y).
\]
D'après 2.5, on a l'isomorphisme:
\[
\underline{\Tor}_i^{\cO_X}(\cO_Y,\cO_Y)
\xrightarrow{\sim}
\bigwedge^i N.
\]

Les \(\underline{\Tor}_i^{\cO_X}(M,\cO_Y)\) sont donc des \(\cO_Y\)-Modules localement libres de type fini, et la classe dans \(K^\circ(Y)\) du complexe \(\mathbb L i^*(\mathbb Ri_*(M))\) est la somme alternée des classes de ses faisceaux d'homologie, soit:
\[
i^*(i_*(x))=\sum_{i\ge0}(-1)^i
\operatorname{cl}\bigl(M\otimes_{\cO_Y}\underline{\Tor}_i^{\cO_X}(\cO_Y,\cO_Y)\bigr)
=\sum_{i\ge0}(-1)^i x.\lambda^i(N).
\]

\textbf{Corollaire 2.8.} Sous les hypothèses de 2.7, soient \(x,y\in K^\circ(Y)\); on a la relation:
\begin{equation}\notag
i_*(xy\lambda_{-1}(N))=i_*(x).i_*(y).
\end{equation}

D'après 2.7, \(y\lambda_{-1}(N)=i^*(i_*(y))\).

Par application de la formule de projection, on trouve donc:
\[
i_*(xy\lambda_{-1}(N))
=i_*(x.i^*(i_*(y)))
=i_*(x).i_*(y).
\]

\section*{3. Calcul du \texorpdfstring{$K^\circ$}{K} d'un schéma éclaté}

\textbf{3.1.} Dans ce paragraphe, on désigne par \(i:Y\to X\) une immersion fermée régulière. On note \(X'\) le schéma obtenu en faisant éclater \(Y\), et \(Y'=Y\times_X X'\).
